\documentclass[../../main.tex]{subfiles}
\begin{document}

{\bf [解]}

\begin{figure}[htbp]
  \centering
  \begin{tikzpicture}[scale=1.25, >=stealth, every node/.style={font=\small}]

    % --- 1. 定数・パラメータ定義 ---
    % a = 0.8, 正方形の一辺は 2
    \def\aVal{0.8}
    \def\sLen{2.0}

    \pgfmathsetmacro{\rOneSq}{1.0 + \aVal*\aVal}
    \pgfmathsetmacro{\rTwoSq}{1.0 + (2.0 - \aVal)*(2.0 - \aVal)}
    \pgfmathsetmacro{\xConnOne}{2.0 + \aVal}      % 2 + a = 2.8
    \pgfmathsetmacro{\xConnTwoSym}{8.0 - \aVal}   % 8 - a = 7.2
    \def\yConn{1.0}

    % --- 2. 曲線 C と x 軸で囲まれた領域のハッチング ---
    \fill[pattern=north east lines, pattern color=blue!40]
    (1, 0) --
    plot[domain=1:\xConnOne, samples=40] (\x, {sqrt(max(0, \rOneSq - (\x - 2.0)^2))}) --
    plot[domain=\xConnOne:5.0, samples=40] (\x, {sqrt(max(0, \rTwoSq - (\x - 4.0)^2))}) --
    plot[domain=5.0:\xConnTwoSym, samples=40] (\x, {sqrt(max(0, \rTwoSq - (\x - 6.0)^2))}) --
    plot[domain=\xConnTwoSym:9.0, samples=40] (\x, {sqrt(max(0, \rOneSq - (\x - 8.0)^2))}) --
    (9, 0) -- cycle;

    % --- 3. 正方形 S の各要所での姿勢（破線表示） ---
    % (1) 初期状態: 底辺 [0, 2], 高さ 2 (P は (1, a))
    \draw[thick, dashed, gray!60] (0, 0) rectangle (2, \sLen);
    \node[gray!70!black, font=\scriptsize] at (0.6, 1.7) {$S$ (初期)};

    % (2) (2,0) まわりに 90° 回転直後: 底辺 [2, 4], 高さ 2 (P は (2+a, 1))
    \draw[thick, dashed, gray!60] (2, 0) rectangle (4, \sLen);

    % (3) (4,0) まわりに 90° 回転直後（対称軸 x=5）: 底辺 [4, 6], 高さ 2 (P は (5, 2-a))
    \draw[thick, dashed, gray!60] (4, 0) rectangle (6, \sLen);

    % (4) (6,0) まわりに 90° 回転直後: 底辺 [6, 8], 高さ 2 (P は (8-a, 1))
    \draw[thick, dashed, gray!60] (6, 0) rectangle (8, \sLen);

    % (5) (8,0) まわりに 90° 回転直後（1周完了）: 底辺 [8, 10], 高さ 2 (P は (9, a))
    \draw[thick, dashed, gray!60] (8, 0) rectangle (10, \sLen);

    % --- 4. 座標軸 ---
    \draw[->, thick] (-0.5, 0) -- (10.6, 0) node[right] {$x$};
    \draw[->, thick] (0, -0.4) -- (0, 2.6) node[above] {$y$};
    \node[below left=2pt] at (0, 0) {$O$};

    % --- 5. 対称軸 x = 5 ---
    \draw[thick, dashed, red!80!black] (5.0, -0.2) -- (5.0, 2.3)
    node[above, font=\footnotesize] {$x = 5$ (対称軸)};

    % --- 6. 曲線 C の各円弧 ---
    \draw[very thick, blue!85!black, domain=1:\xConnOne, samples=50]
    plot (\x, {sqrt(max(0, \rOneSq - (\x - 2.0)^2))});
    \draw[very thick, blue!85!black, domain=\xConnOne:5.0, samples=50]
    plot (\x, {sqrt(max(0, \rTwoSq - (\x - 4.0)^2))});
    \draw[very thick, blue!85!black, domain=5.0:\xConnTwoSym, samples=50]
    plot (\x, {sqrt(max(0, \rTwoSq - (\x - 6.0)^2))});
    \draw[very thick, blue!85!black, domain=\xConnTwoSym:9.0, samples=50]
    plot (\x, {sqrt(max(0, \rOneSq - (\x - 8.0)^2))});
    \node[blue!85!black, font=\normalsize\bfseries] at (2.6, 1.65) {$C$};

    % --- 7. 境界・接続点の補助点線 ---
    \draw[dashed, gray!70] (1, 0) -- (1, \aVal) -- (0, \aVal);
    \node[left=2pt, font=\scriptsize] at (0, \aVal) {$a$};

    \draw[dashed, gray!70] (9, 0) -- (9, \aVal);

    \draw[dotted, gray!80!black, thick] (\xConnOne, 0) -- (\xConnOne, \yConn);
    \draw[dotted, gray!80!black, thick] (\xConnTwoSym, 0) -- (\xConnTwoSym, \yConn);
    \draw[dashed, gray!60] (0, \yConn) -- (10, \yConn);
    \node[left=2pt, font=\scriptsize] at (0, \yConn) {$1$};

    % 回転中心から接続点への半径動径
    \draw[densely dotted, purple!75!black] (2, 0) -- (\xConnOne, \yConn);
    \draw[densely dotted, purple!75!black] (4, 0) -- (\xConnOne, \yConn);

    % --- 8. 主要点プロット ---
    % 回転中心
    \foreach \xc in {2, 4, 6, 8} {
      \fill[black] (\xc, 0) circle (1.2pt);
      \node[below=2pt] at (\xc, 0) {$\xc$};
    }

    % 各姿勢における点 P の位置
    \fill[red!80!black] (1, \aVal) circle (1.4pt) node[above left, font=\scriptsize] {$P$};
    \fill[red!80!black] (\xConnOne, \yConn) circle (1.4pt);
    \fill[red!80!black] (5.0, {2.0 - \aVal}) circle (1.4pt);
    \fill[red!80!black] (\xConnTwoSym, \yConn) circle (1.4pt);
    \fill[red!80!black] (9, \aVal) circle (1.4pt);

    % x 軸上の主要目盛りラベル
    \node[below=2pt] at (1, 0) {$1$};
    \node[below=2pt] at (5, 0) {$5$};
    \node[below=2pt] at (9, 0) {$9$};
    \node[below=3pt, purple!80!black, font=\footnotesize] at (\xConnOne, 0) {$2+a$};
    \node[below=3pt, purple!80!black, font=\footnotesize] at (\xConnTwoSym, 0) {$8-a$};

    % 正方形の一辺の長さガイド（高さ 2）
    \draw[<->, gray!70!black] (-0.2, 0) -- (-0.2, \sLen)
    node[midway, left=1pt, font=\scriptsize] {$2$};

  \end{tikzpicture}
  \caption{正方形 $S$（一辺 $2$）の転がりに伴う各姿勢（点線）と点 $P$ の軌跡 $C$}
  \label{fig:square_rolling_trajectory_with_squares}
\end{figure}

半径$2$の正方形$S$が滑らずにころがるとき，点$P(1,a)$の軌跡$C$を考える．
\begin{enumerate}[(i)]
  \item $S$が点$(2,0)$を中心として時計回りに$90^\circ$回転するとき：
    点$P$は点$(2,0)$を中心とし，半径$\sqrt{(1-2)^2+a^2}=\sqrt{1+a^2}$の円弧を描く．
    $x$座標は$1$から$2+a$まで変化するから，この区間の曲線$C$の方程式は
    \begin{align}
      (x-2)^2+y^2=1+a^2 \quad (1 \leqq x \leqq 2+a) \label{eq:1}
    \end{align}
    である．
  \item $S$の次の頂点が$(4,0)$に接地し，この点を中心として時計回りに$90^\circ$回転するとき：
    点$P$と点$(4,0)$との距離は$\sqrt{(1-2)^2+(a-2)^2}=\sqrt{1+(2-a)^2}$である．
    $x$座標は$2+a$から$5$まで変化するから，この区間の曲線$C$の方程式は
    \begin{align}
      (x-4)^2+y^2=1+(2-a)^2 \quad (2+a \leqq x \leqq 5) \label{eq:2}
    \end{align}
    である．
  \item 同様に，さらに$S$がころがるとき，曲線$C$は直線$x=5$に関して対称である．
    実際，各区間の方程式は
    \begin{align}
      \begin{cases}
        (x-6)^2+y^2=1+(2-a)^2 & (5 \leqq x \leqq 6+a) \\
        (x-8)^2+y^2=1+a^2 & (6+a \leqq x \leqq 9)
      \end{cases} \label{eq:3}
    \end{align}
    となる．
\end{enumerate}

曲線$C$が直線$x=5$に関して対称であることから，求める立体の体積$V(a)$は，$1 \leqq x \leqq 5$の部分の回転体の体積$V_1(a)$の$2$倍に等しい．すなわち
\begin{align}
  V(a) = 2V_1(a) = 2\pi \int_1^5 y^2 \, dx \label{eq:4}
\end{align}
である．そこで以下$V_1(a)$を計算する．
\begin{align}
  V_1(a)
  &= \pi \int_1^{2+a} \left\{(1+a^2)-(x-2)^2\right\} \, dx + \pi \int_{2+a}^5 \left\{(a^2-4a+5)-(x-4)^2\right\} \, dx \\
  &= \pi \int_{-1}^a (1+a^2-x^2) \, dx + \pi \int_{a-2}^1 (a^2-4a+5-x^2) \, dx \label{eq:5}
\end{align}
を各項計算すると
\begin{align}
  (\text{第1項})
  &= \int_{-1}^a (1+a^2-x^2) \, dx \nonumber\\
  &= \left[ (1+a^2)x - \dfrac{x^3}{3} \right]_{-1}^a \nonumber\\
  &= \left( (1+a^2)a-\dfrac{a^3}{3} \right) - \left( -(1+a^2)+\dfrac{1}{3} \right) \nonumber\\
  &= \dfrac{2}{3}a^3 + a^2 + a + \dfrac{2}{3} \label{eq:6}
\end{align}
および
\begin{align}
  (\text{第2項})
  &= \int_{a-2}^1 (a^2-4a+5-x^2) \, dx \nonumber\\
  &= \left[ (a^2-4a+5)x - \dfrac{x^3}{3} \right]_{a-2}^1 \nonumber\\
  &= \left( a^2-4a+5-\dfrac{1}{3} \right) - (a-2)\left( a^2-4a+5 - \dfrac{(a-2)^2}{3} \right) \nonumber\\
  &= a^2-4a+\dfrac{14}{3} - (a-2)\left( \dfrac{2}{3}a^2 - \dfrac{8}{3}a + \dfrac{11}{3} \right) \nonumber\\
  &= a^2-4a+\dfrac{14}{3} - \left( \dfrac{2}{3}a^3 - 4a^2 + 9a - \dfrac{22}{3} \right) \nonumber\\
  &= -\dfrac{2}{3}a^3 + 5a^2 - 13a + 12 \label{eq:7}
\end{align}
である．

したがって\cref{eq:6,eq:7}を\cref{eq:5}に代入して
\begin{align}
  V_1(a) &= \pi \left\{ \left( \dfrac{2}{3}a^3 + a^2 + a + \dfrac{2}{3} \right) + \left( -\dfrac{2}{3}a^3 + 5a^2 - 13a + 12 \right) \right\} \nonumber\\
  &= \pi \left( 6a^2 - 12a + \dfrac{38}{3} \right) \nonumber\\
  &= \pi \left\{ 6(a-1)^2 + \dfrac{20}{3} \right\}
\end{align}
これより，求める体積$V(a)$は
\begin{align}
  V(a) &= 2V_1(a) \nonumber\\
  &= \pi \left( 12a^2 - 24a + \dfrac{76}{3} \right) \nonumber\\
  &= \pi \left\{ 12(a-1)^2 + \dfrac{40}{3} \right\}
\end{align}
である．

$0 \leqq a \leqq 2$において，$(a-1)^2 \geqq 0$であり，等号は$a=1$のとき成り立つ．
よって，$V(a)$が最小となる$a$の値およびその最小値は
\begin{align}
  a=1 \text{ のとき最小値 } \dfrac{40}{3}\pi \quad \cdots\text{（答）}
\end{align}
である．

\end{document}
